\begin{frame}{두 가지 탐험 전략: exploring starts vs
$\varepsilon$-greedy}
  \begin{columns}[T]
    \begin{column}{0.5\textwidth}
      \kb{(1) 탐험적 시작 (exploring starts)}
      \begin{itemize}
        \item 매 에피소드를 \textbf{무작위로 고른} (상태,
          행동) 쌍에서 \textbf{강제로 시작}.
        \item 모든 (상태, 행동) 쌍이 결국 충분히 방문됨.
        \item 이론적으로 깔끔하지만, 임의의 상태·행동에서
          ``시작''할 수 있어야 하므로 \textbf{실전 환경에서는
          항상 가능하지 않다}.
      \end{itemize}
    \end{column}
    \begin{column}{0.5\textwidth}
      \kb{(2) $\varepsilon$-greedy} (Week 2.1 의 그 전략)
      \begin{itemize}
        \item 매 스텝 \textbf{작은 확률로 무작위 행동}을 섞음.
        \item 시작 상태와 무관하게 모든 (상태, 행동) 쌍을
          계속 조금씩 방문.
        \item \textbf{이번 실습은 $\varepsilon$-greedy.}
      \end{itemize}
    \end{column}
  \end{columns}
\end{frame}
